\section{The lists: sources of antisemitic incident data}\label{sec:sources}

Multiple-systems estimation needs $K\ge2$ lists, each an incomplete register of the \emph{same} population $U_d$ of in-scope incidents in a domain $d$. This section maps the candidate US lists onto the canonical design ($k=1$ police, $k=2$ NGO audit, $k=3$ community security organisation, $k=4$ campus or community portal); \cref{tab:src-compare} summarises them. Every figure quoted is attributed to its publisher, none is an input to the simulation study, and published totals from different lists are \emph{not} estimates of a common quantity (\cref{sec:src-harmonise}).

\begin{table}[tbp]
\centering
\caption{Candidate sources for a US application. ``List'' is the role in the canonical $K=4$ design; CST and FRA are international comparators. Entries summarise the sources' own documentation (citations in the text).}\label{tab:src-compare}
\footnotesize
\setlength{\tabcolsep}{3pt}
\renewcommand{\arraystretch}{1.1}
\begin{tabularx}{\textwidth}{@{}>{\raggedright\arraybackslash}p{1.75cm}*{4}{>{\raggedright\arraybackslash}X}@{}}
\toprule
Source (role) & Unit; anti-Jewish identification & Granularity; access & Coverage issues & Dependence risks \\
\midrule
FBI UCR/\allowbreak NIBRS (list 1) & Crimes; NIBRS bias code 21, ``sufficient objective facts'' & Agency $\times$ year tables; records without names or addresses & Voluntary participation; NIBRS-only 2021 & Feeds ADL and Clery; receives community referrals \\
NYPD, California DOJ (list 1, local) & Complaints; events (one or more offences); anti-Jewish category & NYPD rows public by month and precinct; CA annual report & 2023 NYC state publication gap; partial-year CA submissions & Same system as the FBI list \\
ADL Audit, H.E.A.T. Map (list 2) & Incidents, criminal and non-criminal; Israel/Zionism rule since Oct.\ 2023 & Incident listings with redactions & Mid-series scope change & Ingests police and partner reports \\
SCN (list 3) & Threats, suspicious activity; criteria not public & Aggregate summaries; records by agreement & Scope broader than crimes & ADL partner; refers to police \\
CST (UK) & Incidents; 3{,}001 reports screened out (2025) & Annual reports & Police feed disrupted since 2024 & 16\% of 2025 incidents from police \\
Clery data (campus) & Crimes; religion bias only & Institution $\times$ year; public & Superset of target; campus only & Built partly from police data \\
NCVS (survey) & Victimisations; victim perception plus evidence & National, pooled years & Persons and households only & Auxiliary: $\rpol$ \\
FRA (EU) & Respondent-reported incidents & Country; 2018, 2024 waves & Not US & Auxiliary: non-reporting \\
\bottomrule
\end{tabularx}
\end{table}

\subsection{Police-recorded hate crime (list 1)}\label{sec:src-police}

Under the Hate Crime Statistics Act of 1990 \citep{hcsa1990}, the FBI's Uniform Crime Reporting (UCR) Program attached a bias element to the offences already reported to it \citep{fbi2019about}. In the National Incident-Based Reporting System (NIBRS), anti-Jewish bias is code~21, reported only if the investigation reveals sufficient objective facts to lead a reasonable and prudent person to conclude that bias motivated the offender, in whole or in part \citep{fbimanual}. List~1 is therefore a list of \emph{crimes} whose anti-Jewish motivation an agency has \emph{classified}; for 2025 the FBI reported 1{,}528 single-bias anti-Jewish incidents \citep{fbi2026hate,adl2026fbi}. NIBRS holds no names or addresses, so linkage needs the agencies' own records.

\paragraph{Participation and the 2021 transition}
Participation is voluntary at the federal level, and for data year 2021 the FBI accepted only NIBRS submissions \citep{crs2025hate}. The initial 2021 hate-crime release, in December 2022, covered 11{,}834 agencies representing 64.8\% of the US population. After the FBI reopened the collection to summary-format and additional NIBRS reports from 3{,}025 agencies, its supplemental release of March 2023 covered 14{,}859 agencies and 91.1\% of the population \citep{fbi2023supp,doj2023hate2021}. Both formats were accepted again from 2022 \citep{crs2025hate}; for 2024, 16{,}419 participating agencies covered 95.1\% of the population, and the FBI does not estimate counts for non-participants \citep{fbi2025summary}. Participation measures \emph{submission} ($\rsub$ in \cref{sec:funnel}), which can shift abruptly for reasons unrelated to victimisation; the small-area model (\cref{eq:smallarea}) therefore ties police coverage to it.

\paragraph{City and state data}\label{sec:src-local}
City and state data are list~1 observed before federal submission, with richer linkage keys. The NYPD Hate Crimes dataset has one public row per confirmed complaint, with month, precinct, offence and a bias value \texttt{ANTI-JEWISH} \citep{nypdhate}, which makes New York City a natural pilot (\cref{sec:plan}); the state published no 2023 data for the city because of the NYPD's NIBRS transition \citep{osc2024hate}. California counts hate-crime \emph{events} under Penal Code sections 13023 and 422.55, with 289 anti-Jewish events for 2023, 310 for 2024 and 289 for 2025, not all from full-year submissions \citep{cadoj2025hate,cadoj2026hate}; such units need not match federal incidents.

\subsection{The NGO audit (list 2)}\label{sec:src-audit}

ADL's Audit, kept since 1979, counts criminal and non-criminal incidents of assault, harassment and vandalism, assesses each report for credibility, counts each incident once, and can treat a related series of acts as one incident \citep{adl2025audit,adl2026audit}. Since 7 October 2023 it has included certain expressions of opposition to Zionism that could be perceived as supporting terrorism or attacks on Jews, Israelis or Zionists, while excluding legitimate political protest \citep{adl2025audit,adl2026audit}. ADL reported that 58\% of its 2024 incidents contained elements related to Israel or Zionism, and recorded 9{,}354 incidents for 2024, the highest in the Audit's history, and 6{,}274 for 2025 \citep{adl2025audit,adl2026audit}. Totals before and after October 2023 therefore measure partly different populations.

The Audit is a \emph{compiled} list: its largest source is victims' own reports \citep{adl2025audit}, but from 2021 it incorporated reports from partner Jewish organisations, among them the Secure Community Network and Hillel International; ADL attributed 494 incidents in 2021 to these partnerships \citep{adl2022audit}. Some of its records thus originate in lists~1 and~3. The H.E.A.T. Map publishes redacted incident-level listings \citep{adlheatmap}, enough for coarse linkage on date, place and category. Provenance (how ADL learned of each incident) is not, to our knowledge, public; it is the most valuable field to request, because it turns referral into an identification strategy (\cref{prop:referral}(c)--(d)).

\subsection{Community security organisations and campus data (lists 3 and 4)}\label{sec:src-community}

The Secure Community Network (SCN), founded in 2004, describes itself as the official safety and security organisation of the Jewish community in North America \citep{scnwho}. Its Duty Desk shares threat and incident information with law enforcement and partners \citep{scnmonitoring}, and it publishes only aggregate summaries \citep{scn2024monthly}, so its use as a list requires an agreement and a mapping of its scope (threats, suspicious activity) onto the common ontology; as an Audit partner that refers to the police, it is linked to both other lists. The UK's Community Security Trust (CST) shows what a mature community list documents. It recorded 3{,}700 antisemitic incidents in 2025, up from 3{,}556 in 2024; it excluded a further 3{,}001 reports that it did not deem antisemitic; 587 (16\%) of its 2025 incidents came from police forces under information-sharing agreements; and since March 2024 a technical problem has disrupted the flow of reports from the Metropolitan Police \citep{cst2026incidents}. A US community list needs the same documentation of screening, provenance and feed interruptions.

\paragraph{Campus data}\label{sec:src-campus}
The Clery Act \citep{cleryact} and 34~C.F.R.\ \S~668.46 \citep{ecfr668} require institutions in federal student-aid programmes to disclose hate crimes in their Clery geography, published at institution level \citep{edcss}. The bias categories include religion but not anti-Jewish bias, which institutions were invited in 2024 to disclose only voluntarily and outside the statistical grids \citep{ed2024subcat}. Clery data thus cover a \emph{superset} of the target, harmful in a capture list (\cref{ex:src-scope}), and are partly built from police statistics \citep{ecfr668}. For list~4 we therefore propose a campus or community self-report portal that identifies anti-Jewish bias, used under agreement and only within the campus stratum, where its remit is declared (\cref{asm:fw-scope}(b), \cref{sec:plan-models}).

\subsection{Victimisation surveys: anchors, not lists}\label{sec:src-surveys}

The National Crime Victimization Survey (NCVS) counts a victimisation as a hate crime when the victim perceives bias and reports hate language, hate symbols or police confirmation, whether or not it was reported to police, and the Bureau of Justice Statistics publishes the share reported \citep{bjs2021hate}; others use these data to study who reports \citep{zaykowski2010,lantz2022,pezzella2019}. It misses institutional victims and has, to our knowledge, no estimate for anti-Jewish bias, so it serves only for triangulation (\cref{sec:plan-validation}). The EU Agency for Fundamental Rights has surveyed Jewish people directly: 16{,}395 respondents in 12 Member States in 2018, of whom 79\% of those who had experienced antisemitic harassment had not reported the most serious incident to the police or to any other organisation \citep{fra2018antisemitism}, and almost 8{,}000 respondents in 13 Member States in a third survey published in 2024 \citep{fra2024antisemitism}. We do not transport these numbers to the United States, but non-reporting to \emph{any} organisation means that many incidents never reach a recording channel; they lie outside $N_V$ and are invisible to every list (\cref{sec:funnel}).

\subsection{Definitional harmonisation}\label{sec:src-harmonise}

The lists count different things (crimes, incidents, threats) in different units, so we code every record on five axes: (i) \emph{conduct} (assault; threat or intimidation; property damage; other offence; non-criminal conduct); (ii) \emph{bias} (anti-Jewish as the single bias, among multiple biases, or not established); (iii) \emph{target} (individual, institution, public space); (iv) \emph{setting} (campus, online, other); and (v) a \emph{content flag} for Israel- or Zionism-related expression. The primary estimand (\cref{sec:plan}) counts \emph{criminal antisemitic incidents}: conduct that is an offence under the jurisdiction's criminal law, with anti-Jewish bias meeting a stated evidentiary standard. Capture follows each list's own classification: $Y_{ik}=1$ if list $k$ holds a record of $i$ that it classifies as antisemitic. An antisemitic crime that police recorded without a bias code is therefore \emph{not} on list~1; that loss is the $\rcls$ stage of the funnel, not a linkage failure. Scope is coded per record, blind to linkage (otherwise incidents on several lists would have more chances to be coded in scope), by a rule fixed in advance and varied in sensitivity analyses.

Every capture--recapture model assumes that each list \emph{could} capture every member of $U_d$ (\cref{asm:fw-scope}). If list $k$ cannot capture a sub-population, for example the police list and non-criminal harassment, then $\Pr(Y_{ik}=1)=0$ there, and this \emph{structural zero} masquerades as heterogeneity or list interaction.

\begin{example}[Scope mismatch as hidden heterogeneity]\label{ex:src-scope}
Let $U=A\cup B$, where $A$ holds $N_A$ criminal incidents and $B$ holds $N_B$ non-criminal ones, and suppose the target is $N=N_A+N_B$. List~1 captures only incidents in $A$, each with probability $p_1$. List~2 captures incidents independently of list~1, with probability $p_{2A}$ in $A$ and $p_{2B}$ in $B$. Then $\E n_{1+}=N_Ap_1$, $\E n_{2+}=N_Ap_{2A}+N_Bp_{2B}$ and $\E n_{11}=N_Ap_1p_{2A}$, so the Lincoln--Petersen estimator converges to
\begin{equation}\label{eq:src-scope}
  \frac{\E n_{1+}\,\E n_{2+}}{\E n_{11}} \;=\; N_A + N_B\,\frac{p_{2B}}{p_{2A}}.
\end{equation}
This equals $N$ only if $p_{2B}=p_{2A}$ and equals $N_A$ only if $p_{2B}=0$; in general it estimates neither. With $K>2$ lists the same mechanism appears as spurious interactions and contaminates the extrapolation to $\nzero$.
\end{example}

There are two remedies: restrict every list to the common scope, so that the estimand is $N_A$ (our primary choice), or model the strata separately with declared structural zeros, which requires at least two lists that can capture $B$. The Audit's October 2023 scope change is a structural zero in \emph{time}, so domains should not straddle it. The lists differ most on Israel- or Zionism-related expression, as do the IHRA working definition \citep{ihra2016} and the Jerusalem Declaration on Antisemitism \citep{jda2021}; we do not adjudicate, but report the estimand with and without flagged incidents. For the primary, crimes-only estimand the flag matters only where an offence is also present.

\subsection{Dependence between lists}\label{sec:src-dependence}

Four mechanisms make the lists dependent. (a) \emph{Referral}: police and partners feed the Audit, police feed Clery, and community organisations feed the police; \cref{prop:referral} models it as forwarding a fraction $q$ of one list's captures to another. (b) \emph{Heterogeneity}: severity, visibility and media attention raise capture on all lists at once; ignored, this biases Lincoln--Petersen downward for two lists, while for $K\ge3$ either sign can occur (\cref{prop:id-hetero,prop:direction}). (c) \emph{Specialisation}: lists designed to reach what others miss, such as a campus portal, produce negative dependence, which biases two-list estimates upward; for $K\ge3$ the direction again follows the oriented top interaction. (d) \emph{Unidentified residual dependence}: $\omega$ is not identified (\cref{thm:nonid}) and is treated as a sensitivity parameter (\cref{thm:closed,cor:bounds}). Linkage error interacts with all four, distorting the table before any dependence model is fitted (\cref{thm:linkbias}); its uncertainty is propagated by mixing over linkage draws (\cref{alg:pipeline}).
